\section{2026 竞争格局：模型能力已进入“多维竞争”}

\subsection{DeepSeek 时刻与中美开源策略分化}

对话把 2025 年初 DeepSeek R1 的影响称作一次“加速度事件”。它的意义不只是某个榜单名次，而是证明高质量开源权重模型可以快速重置市场预期：开发者会重新评估 API 依赖、企业会重新评估私有部署路径、研究者会重新评估可复现基线。

Nathan 特别提到中国团队并非单点突破，而是出现了多点并进：Z.AI、Minimax、Moonshot 等都在同一时间窗口内交出可用模型。这个现象背后不是“某条秘密算法”，而是工程组织在高压竞争中的并行推进能力。Sebastian 的判断也很直接：\textit{``ideas are not proprietary; resources are''}。换句话说，差异常在执行资源，而不在概念独占。

\begin{importantbox}{开源许可正在成为商业武器}
过去讨论开源，常把焦点放在“技术透明度”；现在更现实的焦点是“市场进入方式”。
\begin{itemize}[leftmargin=1.5em]
    \item 限制条款较少的开源权重，可快速扩大开发者触达；
    \item 对企业而言，可控部署比“最佳 benchmark”更重要；
    \item 对模型提供方而言，开放本身可以换来生态心智与后续商业化筹码。
\end{itemize}
\end{importantbox}

\begin{warningbox}{“开源=免费午餐”是误解}
开源权重降低了接入成本，但没有消除工程成本。企业依然要支付推理资源、数据清洗、评测体系、日志治理、合规审计与模型回滚的维护成本。真正难的是把模型能力转成稳定服务，而不是把 checkpoint 下载下来。
\end{warningbox}

\subsection{头部闭源模型的真实优势}

访谈里对闭源阵营的判断很克制。OpenAI 的优势被描述为“持续落地新研究方向并快速产品化”；Google 的优势被描述为“垂直整合的算力与基础设施”；Anthropic 的优势被描述为“在 coding 场景里形成了强产品共识”。这三者本质上是不同类型的组织能力。

\begin{knowledgebox}{为什么同一时期会出现“认知错位”}
社交媒体讨论热度和真实用户规模经常错位：
\begin{itemize}[leftmargin=1.5em]
    \item 一些模型在 X 上热度很高，但并不代表企业渗透率高；
    \item 一些模型在开发者圈口碑一般，但在大众入口里增长很快；
    \item 真正决定长期份额的，是默认入口、价格结构、延迟体验和稳定性，而不只是单次榜单排名。
\end{itemize}
\end{knowledgebox}

\subsection{本章小结}

2026 的竞争已经不是“闭源 vs 开源”的二元战，而是许可、成本、入口、工具、品牌五个维度的组合博弈。DeepSeek 时刻只是把这个趋势加速显化。

